\begin{helper}
Let \(N,n,m\) be natural numbers with \(0<N\), \(n\le N\), and
\(m\le N\), and let \(M\subseteq \operatorname{Fin}(N)\) have cardinality
\(m\).  Put \(q=m/N\).  Let \(v_i=1\) if \(i\in M\), and \(v_i=0\)
otherwise.  For \(X\in \operatorname{Fin}(N)^n\), let
\[
c_i(X)=|\{j\in \operatorname{Fin}(n):X_j=i\}|.
\]
Then for every \(L\in\mathbb R\),
\[
(1-q+q e^L)^n
=\frac{1}{N^n}\sum_{X\in \operatorname{Fin}(N)^n}
  \frac{1}{N!}\sum_{\pi\in S_N}
  \exp\!\left(L\sum_{i=1}^N c_i(X)v_{\pi(i)}\right).
\]
\end{helper}
\begin{proof}
Because \(0<N\), the finite set \(U=\operatorname{Fin}(N)\) has
cardinality \(N\).  For a fixed permutation \(\pi\in S_N\) and a fixed
sequence \(X=(X_1,\ldots,X_n)\in U^n\), the definition of the count vector
gives the exact identity
\[
\sum_{i=1}^N c_i(X)v_{\pi(i)}
 =\sum_{j=1}^n v_{\pi(X_j)}.
\]
Indeed, each occurrence of a value \(i\) among the coordinates of \(X\)
contributes one copy of \(v_{\pi(i)}\) to the right hand side, and the
number of such occurrences is precisely \(c_i(X)\).

Therefore, for fixed \(\pi\),
\[
\frac{1}{N^n}\sum_{X\in U^n}
  \exp\!\left(L\sum_i c_i(X)v_{\pi(i)}\right)
=
\frac{1}{N^n}\sum_{X\in U^n}
  \prod_{j=1}^n \exp\!\left(L\,v_{\pi(X_j)}\right).
\]
The sum over \(U^n\) factors into the product of \(n\) identical one-coordinate
sums:
\[
\frac{1}{N^n}\sum_{X\in U^n}\prod_{j=1}^n
\exp\!\left(L\,v_{\pi(X_j)}\right)
=
\left(\frac1N\sum_{a\in U}\exp(L\,v_{\pi(a)})\right)^n .
\]
Since \(\pi\) is a bijection of \(U\),
\(\sum_{a\in U}\exp(L\,v_{\pi(a)})=\sum_{a\in U}\exp(L\,v_a)\).
Exactly \(m\) of the \(v_a\)'s are equal to \(1\), because
\(|M|=m\), and the remaining \(N-m\) are equal to \(0\).  Hence
\[
\frac1N\sum_{a\in U}\exp(L\,v_a)
=\frac{N-m}{N}\exp(0)+\frac{m}{N}\exp(L)
=1-q+q e^L .
\]
Thus the inner average over \(X\), for every fixed \(\pi\), is
\((1-q+q e^L)^n\).

The outer average over permutations has total weight
\((N!)^{-1}\cdot |S_N|=(N!)^{-1}N!=1\).  Averaging the same value
\((1-q+q e^L)^n\) over all \(\pi\in S_N\) therefore leaves it unchanged.
This is exactly the displayed identity.
\end{proof}
